\chapter{Discusiones} \label{discusiones}

\section{Estudio transversal sobre la influencia de marcadores en el declive cognitivo en la Enfermedad de Parkinson}
Los resultados del análisis en Python revelan la influencia de las variables estudiadas en el deterioro cognitivo de pacientes con Enfermedad de Parkinson.\\

Observando las dos tablas resultantes, se puede llevar a algunas conclusiones sobre los distintos parámetros. En la tabla \ref{tab:tabla3}, tras la comparación de los valores en la linea de base, se observan diferencias relevantes en algunas filas entre el grupo de pacientes con Parkinson y el grupo de control. En primer lugar, cabe destacar que para que el estudio sea fiable, ambos grupos tiene las mismas características demográficas. Sin embargo, entre los parámetros cognitivos y clínicos se encuentran varios marcadores con indicador de significancia en la columna del Post-hoc. Entre los que se pueden señalar están GDS, REM, MOCA, SDMT y MEANPUT. \\

Por otra parte, en la tabla \ref{tab:tabla2} destaca la relación entre las distintas características de estudio y el desarrollo de diferentes niveles de declive cognitivo en los pacientes PD. Los valores de p-valor y los indicadores de post-hoc señalan que variables pueden tener un efecto significativo en el transcurso del deterioro cognitivo en los distintos subgrupos. Entre ellas, diferenciadas por tener un \textit{valor p} igual o por debajo de 0.001, se ve como por ejemplo el SDMT y el MEANCAU son considerables tanto en la comparación entre sujetos cognitivamente normales y con deterioro cognitivo como entre cognitivamente normales y los pacientes que desarrollaron demencia. Sin embargo, hay factores que solo muestran efectos significativos en una de las dos comparativas, marcados con un solo indicador (\textit{a} o \textit{b}). Entre estos últimos se encuentran REM, con diferencias importantes solo entre PD\_CU y PD\_MCI; o A$\beta_{42}$, con diferencia notable entre PD\_CU y PDD. En esta tabla, dentro de las variables demográficas, la edad es un indicador del avance de los pacientes hacia un deterioro cognitivo medio y de la aparición de la demencia. Además, se puede ver como las variables cognitivas tienen una gran importancia en el desarrollo del declive cognitivo, sobre todo indicando la evolución de PD\_CU a PD\_MCI. Aún así, los parámetros LNS, SDMT, INMEDIATE y SFT también marcan la evolución de los sujetos hacia el grupo con demencia. Entre las variables clínicas, ABETA y PTAU señalan la evolución hacia el colectivo PDD, y MEANCAU es una característica significativa tanto de la aparición del deterioro cognitivo como de la demencia. \\

Es importante destacar las variables que resultaron significativas en el análisis, tanto en la tabla \ref{tab:tabla3} como en la tabla \ref{tab:tabla2}. Todas ellas son medidas cognitivas y clínicas. Entre estas, se encuentran: UPSIT, SDMT, INMEDIATE, SFT y MEANCAU. Estos parámetros pueden tener implicaciones cruciales para el diagnostico anticipado y el desarrollo de estrategias para el tratamiento.\\

En resumen, en los resultados de este estudio destacan marcadores significativos en el declive cognitivo en la Enfermedad de Parkinson. Estos hallazgos tienen el potencial de mejorar la predicción temprana de la enfermedad y permitir intervenciones más efectivas. La identificación de estos marcadores puede conducir a un enfoque más personalizado en el tratamiento de pacientes con Parkinson, mejorando así su calidad de vida y ofreciendo esperanza a quienes se ven afectados por esta enfermedad. También es importante tener en cuenta que los resultados obtenidos pueden variar según la población y otros factores.


\section{Análisis longitudinal y reconstrucción temporal de un único sujeto}

En las gráficas que podemos ver en las figuras \ref{fig:conv1_gvsource}, \ref{fig:conv1_tauxi} y \ref{fig:conv1_noise} se demuestra la convergencia de los parámetros estimados por Leaspy a lo largo de 50 iteraciones, indicando la validez de los parámetros (g, v0, sources, tau, xi, noise).\\

En el caso de las gráficas obtenidas que estiman las trayectorias \textit{'short-term'} de cada variable, se puede observar en la figura \ref{fig:trayM5norm} como en los parámetros ABETA, MEANCAU Y HVLTRDLY su estimación no se cruza con la curva de unión de los puntos de sus observaciones. En las otras tres (SDMT, LNS, SFT) si que se puede ver una clara intersección entre las rectas, ya que a mayor numero de mediciones mas precisa es la estimación.


\section{Modelos de progresión con múltiples sujetos}
Comparando de forma general los dos métodos de modelado utilizados en este trabajo, a pesar de que Leaspy obtiene resultados robustos para la predicción de enfermedades como el Alzheimer, muestra ciertas variaciones para el caso del Parkinson. Además, en la mayoría de las métricas evaluadas, RPDPM (\textit{Robust Parametric Disease Progression Modeling}) tiene un rendimiento superior para el conjunto de todos los modelos estimados.\\

Para el caso de la detección de pacientes PD\_CU (cognitivamente normales) en el grupo de \textit{test}, se detectan valores similares para ambos métodos, siendo ligeramente superiores para RPDPM, sobre un 0,1\% mayores que los de Leaspy. Para el caso de la detección de sujetos PD\_MCI (deterioro cognitivo leve), hay una diferencia mayor, siendo entorno a un 20\% más exacto utilizando RPDPM. En la correlación con la reserva cognitiva Leaspy da una media superior a RPDPM, sobre un 2\% por encima, mientras que en el Área Bajo la Curva (AUC), RPDPM vuelve a mostrar valores superiores a Leaspy, siendo un 5\% más preciso.\\

Centrando la discusión en la aplicación clínica, para el caso de RPDPM \cite{ghazi2021robust} los valores son bastante similares para los vectores de menor dimensión formados por tres biomarcadores y para los de mayor dimensión, compuestos por cuatro o cinco características. La variación en todas las prestaciones es menor del 0,3\%. En el caso de Leaspy, hay una diferencia más grande entre los vectores de mayor y menor dimensión, obteniendo mejores resultados en el caso de los de tamaño mayor. Esto indica por un lado, la posibilidad de utilizar el método de RPDPM tomando un menor número de muestras en los pacientes, ya que harían falta observaciones de un menor número de marcadores para obtener resultados precisos. En el caso de Leaspy, las predicciones llevadas a cabo con un menor numero de características son menos precisas que las que emplean más marcadores para el mismo método. También muestran peores resultados que los obtenidos para el mismo número de biomarcadores con el método RPDPM. En definitva, para obtener resultados óptimos con Leaspy, es preferible utilizar un conjunto de características más amplio. A diferencia de RPDPM, donde la reducción en la cantidad de biomarcadores no afecta significativamente la precisión, Leaspy parece depender en mayor medida de un conjunto más completo de datos para mantener su robustez en las predicciones. \\

La combinación de características que mejores resultados ha dado para los modelos generados en RPDPM es upsit, LNS, SDMT y ADL y para Leaspy es updrs\_totscore, tau y asyn. Los valores del porcentaje de pacientes pCU detectados son respectivamente 97.46\% y 100\%, el porcentaje de pacientes pMCI es 32.47\% y 6.96\%, y la correlación con la reserva 65.41\% y 100\% para cada caso. En los valores medios de todas las combinaciones generadas, se ha visto que Leaspy muestra porcentajes más bajos. No obstante, en estos valores se refleja como, a excepción de la medida de la detección de pacientes con deterioro cognitivo leve,  en los otros dos parámetros, Leaspy muestra valores más altos. Esto refleja una mayor variabilidad en las predicciones de los modelos evaluados con Leaspy.\\

Por otro lado, el vector de RPDPM formado por los tres biomarcadores sdmt, bjlot y hvltrdly da mejores resultados que la combinación de mayor dimensión para el caso del conjunto de pacientes de prueba, siendo un factor positivo para la aplicación clínica de este método. 

En cuanto al desplazamiento espacial de las trayectorias pronosticadas, medido a través del Error Absoluto Medio (MAE), se observa como para aquellos modelos en los que se obtuvieron los mejores resultados de diagnóstico clínico, también muestran valores mas bajos en el MAE de los biomarcadores característicos. Esto indica una mayor precisión en la predicción y estimación de las curvas de evolución de la enfermedad, lo cual sucede tanto para RPDPM como para Leaspy.


